\documentclass[12pt]{article}
\usepackage{circuitikz}
\usepackage{pgfplots}
\usepackage{amsmath}
\usepackage{siunitx}

\begin{document}

\section*{Circuit RC}

\subsection*{Schéma}
\begin{center}
\begin{tikzpicture}[american]
    % Composants
    \draw (0,0) node[ground] {} to[V=\SI{10}{\volt}] ++(0,3)
          to[R=\SI{1k}{\ohm}] ++(3,0)
          to[C=\SI{1\mu}{\farad}] ++(0,-3) -- (0,0);

    % Étiquettes
    \node at (0+1.5,3) [above] {$R$};
    \node at (0+3,3) [above] {$C$};
\end{tikzpicture}
\end{center}

\subsection*{Graphique Théorique}
La tension suit la loi :
\[ V_C(t) = V_S(1-e^{-t/RC}) \]

\begin{center}
\begin{tikzpicture}
\begin{axis}[
    xlabel={Temps ($\si{\second}$)},
    ymin=0, ymax=10,
    samples=100
]
\pgfmathsetmacro{\RC}{1000*1e-6}
\addplot[blue, domain=0:0.005] {10*(1-exp(-x/\RC))};
\end{axis}
\end{tikzpicture}
\end{center}

\end{document}